\section{Discussion}
\label{sec:discussion}

\subsection{Traditional models and deep learning}
The latest notebook execution supports a restrained conclusion: traditional machine learning and deep learning perform similarly on these tabular survey indicators. E07 tuned gradient boosting leads the validation selection at macro AP 0.4812, while E17 Functional wide-and-deep is the best deep model at 0.4799. On test data their macro ROC-AUC values are 0.8011 and 0.8004, and macro AP values are 0.4792 and 0.4779. This is a small separation, not evidence that one paradigm dominates. Logistic regression remains a useful transparent baseline (test macro AP 0.4690), close enough to make complexity and interpretability part of the practical decision.

The experiment sequence also suggests that adding capacity is not a reliable route to improvement. The two feature-engineering comparisons move macro AP only marginally, and the deep model with engineered features improves over E16 by about 0.0002 on validation AP. The dataset's indicators encode broad health status, age and healthcare contact, so flexible networks may have little new signal to discover. The search results favour regularisation and modest tree complexity rather than ever larger models. These observations are consistent with tabular-learning literature, but they should be read as evidence about this dataset and split rather than a general rule for all health data.

\subsection{Threshold choice and screening burden}
A ranking metric does not determine how many people should be referred for follow-up. At the default threshold, the E07 classifier catches only 16.8\% of diabetes cases. Lowering the threshold to 0.117 raises sensitivity to 84.0\%, but specificity falls to 68.8\%, PPV to 28.7\%, and the flag rate rises to 38.1\%. This is a sizeable screening workload. The threshold was selected by validation F2, which values recall more than precision; it is not a medical recommendation and should not be transferred to a clinic without local evaluation, calibration and a defined follow-up capacity.

The class-weighted experiments and validation-tuned decision weights also illustrate that class recall can move substantially without a corresponding gain in threshold-free ranking. A report that emphasised accuracy alone would conceal the fact that prediabetes is rarely recognised and that a high-sensitivity diabetes rule generates many false-positive flags. The useful output for a screening workflow would be a prioritisation for confirmatory testing, not a diagnosis.

\subsection{Labels, explainability and equity}
The outcome records whether a respondent was told they had diabetes. It cannot distinguish a true negative from a person who has diabetes but has not been tested or informed. As a result, some false-positive flags may resemble plausible undiagnosed cases, while people with limited access to care can be mislabelled in the training target. The weak prediabetes performance is especially consequential because the category is both uncommon and dependent on prior testing. The current models learn the survey's diagnosis process as well as underlying health patterns.

The explanation outputs identify BMI, general health, age, hypertension and cholesterol as useful signals, but these features should not be interpreted as causes. Check-up recency and cholesterol testing partly represent contact with the health system; a diagnosis can change both behaviour and access to care. Attribution methods also distribute importance differently across raw and engineered features, so no single ranking is definitive. Excluding race from model inputs does not eliminate subgroup performance differences. The current subgroup outputs show low recall among younger adults and variation across race/ethnicity and income, which warrants careful local assessment rather than a claim that the model is fair or unfair.

\subsection{Implications for rural sub-Saharan Africa}
The project is motivated by communities where laboratory screening may be difficult to access, but BRFSS is a US telephone survey from 2015. Its predictors, diagnosis pathways, and prevalence are not a substitute for data collected in rural African settings. The results support investigating low-cost questionnaire-based triage, not deploying this model directly. Logistic regression is attractive as a transparent, lightweight baseline; however, local labels and clinical validation matter more than choosing a sophisticated architecture.

A responsible next step would collect laboratory-confirmed glycaemic measurements alongside questionnaire indicators, include communities with varied access to care, and evaluate sensitivity, calibration and referral burden by subgroup. Decision thresholds should be agreed with health workers and clinics based on the availability of confirmatory tests. Only after external and prospective validation should a model be considered for real screening support.

\subsection{Limitations and next work}
This analysis uses one cross-sectional year, self-reported labels, and a subsample for hyperparameter search. The test metrics are not a complete design-based evaluation of BRFSS's complex survey structure, and calibration was not comprehensively assessed. Results may shift across years or populations. Future work should compare the classifier with laboratory-defined outcomes, use survey-aware estimation, evaluate calibration and threshold stability, and validate on later BRFSS data and locally collected African cohorts. These steps would address the main limitations before any clinical use.
